% Compile with XeLaTeX, TeXLive 2013 or more recent
\documentclass{beamer}

% Base packages
\usepackage{fontspec}
\defaultfontfeatures{Scale=MatchLowercase, Mapping=tex-text}

\usepackage{amsfonts}
\usepackage{amsmath}
\usepackage{longtable}
\usepackage{csquotes}
\usepackage{standalone}

\usepackage{graphicx}
\graphicspath{{./images/}}

\usepackage{tikz}
\usetikzlibrary{shapes, calc, arrows, decorations, patterns, chains, fit, petri}
\usetikzlibrary{backgrounds, positioning, decorations.pathreplacing}

\newcommand{\inputpicture}[1]{\input{../pictures/#1}}

\usepackage{tabularx}
\usepackage{multirow}

\usepackage{listings}
\lstset{language=C, basicstyle=\ttfamily, breaklines=true, keepspaces=true, keywordstyle=\color{blue}}

% Setup fonts
\newfontfamily\russianfont{CMU Serif}
\setromanfont{CMU Serif}
\setsansfont{CMU Sans Serif}
%\setmonofont{CMU Typewriter Text}

% Be able to insert hyperlinks
\usepackage{hyperref}
\hypersetup{colorlinks=true, linkcolor=black, filecolor=black, citecolor=black, urlcolor=blue , pdfauthor=Evgeny Yulyugin <yulyugin@gmail.com>, pdftitle=Параллельное программирование}
% \usepackage{url}

% Misc optional packages
\usepackage{underscore}
\usepackage{amsthm}

% A new command to mark not done places
\newcommand{\todo}[1][]{{\color{red}TODO\ #1}}

\newcommand{\abbr}{\textit{англ.}\ }

\newif\ifmipt
\newif\ifsbertech
\input{target}

\ifmipt
\subtitle{Курс <<Параллельное программирование>>}
\fi
\ifsbertech
\subtitle{Курс <<Инфраструктура многопроцессорных систем>>}
\fi
\subject{Лекция}
\author[Евгений Юлюгин]{Евгений Юлюгин \texorpdfstring{\\}{Lg} \small{\href{mailto:yulyugin@gmail.com}{yulyugin@gmail.com}}}
\date{\today}
\pgfdeclareimage[height=0.5cm]{mipt-logo}{common/mipt.png}
\logo{\pgfuseimage{mipt-logo}}

\typeout{Copyright 2014 Evgeny Yulyugin}

\usetheme{Berlin}
\setbeamertemplate{navigation symbols}{}%remove navigation symbols


\title{Технология MPI. Модель MPI приложения. Базовые функции MPI.}

\begin{document}

\begin{frame}
\titlepage
\end{frame}

\section*{Обзор}

\begin{frame}{На этой лекции}
\tableofcontents
\end{frame}

\section{История MPI}

\begin{frame}{Интерфейс передачи сообщений}
Интерфейс передачи сообщений (\abbr Massage Passing Interface, \texttt{MPI}) --- программный интерфейс, позволяющий обмениваться сообщениями между процессами выполняющими одну задачу.
\vfill
\begin{itemize}
    \item Разработан Уильямом Гроуппом, Эвином Ласком и другими;
    \item Первая версия разрабатывалась в 1993-1994 гг;
    \item \texttt{MPI} 1.1 опубликован в 12 июня 1995 года. Поддерживается
    большинством современных реализаций \texttt{MPI}. Первая раелизация
    появилась в 2002 году;
    \item Текущая версия --- \texttt{MPI} 3.0. 2012~г.;
    \item Существуют реализации для языков Fortran, Java, C и C++.
\end{itemize}
\end{frame}

\begin{frame}{Цель MPI}
Основная цель:
\begin{itemize}
    \item Стандартизация,
    \item Переносимость кода,
    \item Эффективная реализация,
    \item Богатство возможностей,
    \item Поддержка архитектур с неоднородной памятью.
\end{itemize}
\end{frame}

\section{Базовые функции MPI}

\begin{frame}[fragile]{Инициализация и завершение процессов}
Определены в заголовочном файле \texttt{mpi.h}
\vfill
C:
\begin{lstlisting}
int MPI_Init(int *argc, char ***argv);
int MPI_Finalize(void);
\end{lstlisting}
\vfill
C++:
\begin{lstlisting}
void MPI::Init(int& argc, char**& argv);
void MPI::Init();
void MPI::Finalize();
\end{lstlisting}
\end{frame}

\begin{frame}[fragile]{MPI <<Hello world!>>}
\begin{lstlisting}
#include <stdio.h>
#include <mpi.h>

int main (int argc, char **argv) {
    MPI_Init(&argc, &argv);
    printf("Hello World!\n");
    MPI_Finalize();
    return 0;
}
\end{lstlisting}
\end{frame}

\begin{frame}{Коммуникатор MPI}
\begin{itemize}
    \item Группа процессов, которые могут взаимодействовать друг с другом,
    \item Все обращения к MPI функциям содержат коммуникатор, как параметр,
    \vfill
    \item Виды коммуникаторов:
    \begin{itemize}
        \item \texttt{MPI_COMM_WORLD} --- Создается при вызове \texttt{MPI_Init}.
        Содержит все процессы программы.
        \item \texttt{MPI_COMM_SELF} --- Содержит только сам процесс,
        \item Пользовательские коммуникаторы.
    \end{itemize}
\end{itemize}
\end{frame}

\begin{frame}[fragile]{Self-identification}
Размер коммуникатора.
\begin{lstlisting}
int MPI_Comm_size(MPI_Comm comm, int *size);
\end{lstlisting}
\vfill
Rank --- идентификационный номер процесса.
\begin{itemize}
    \item Начинается с 0 до (n-1), где n -- число процессов,
    \item Используется для определения процесса-получателя/отправителя.
\end{itemize}
\begin{lstlisting}
int MPI_Comm_rank(MPI_Comm comm, int *rank);
\end{lstlisting}
\end{frame}

\begin{frame}[fragile]{Шаблон}
\begin{lstlisting}[basicstyle=\scriptsize]
#include <stdio.h>
#include <mpi.h>

int main (int argc, char **argv) {
    int rank, size;

    /* Non-parallel part of code here. */

    MPI_Init(&argc, &argv);
    MPI_Comm_size(MPI_COMM_WORLD, &size);
    MPI_Comm_rank(MPI_COMM_WORLD, &rank);

    /* Parallel code here. */

    MPI_Finalize();
    return 0;
}
\end{lstlisting}
\end{frame}

\begin{frame}[fragile]{Hello, MPI World!}
\begin{lstlisting}[basicstyle=\scriptsize]
#include <stdio.h>
#include <mpi.h>

int main (int argc, char **argv) {
    int rank, size;
    MPI_Init(&argc, &argv);
    MPI_Comm_size(MPI_COMM_WORLD, &size);
    MPI_Comm_rank(MPI_COMM_WORLD, &rank);
    printf("Hello MPI World! I'm %d of %d.\n", rank, size);
    MPI_Finalize();
    return 0;
}
\end{lstlisting}
\end{frame}

\section{Передача сообщений}

\begin{frame}{Сообщения}
Сообщение --- массив элементов заданного типа данных MPI.
\vfill
Типы данных в MPI:
\begin{itemize}
    \item Базовые типы,
    \item Производные типы.
\end{itemize}
\vfill
Производные типы строятся с использованием базовых типов.
\end{frame}

\begin{frame}{MPI_Datatype}
\begin{table}[htpb]
    \begin{center}
    \begin{tabular}{|l|l|}
    \hline
    MPI_Datatype        &   C Datatype      \\
    \hline
    MPI_CHAR            &   Signed char     \\
    MPI_SHORT           &   Signed short    \\
    MPI_INT             &   Signed int      \\
    MPI_LONG            &   Signed long     \\
    MPI_UNSIGNED_CHAR   &   Unsigned char   \\
    MPI_UNSIGNED_SHORT  &   Unsigned short  \\
    MPI_UNSIGNED        &   Unsigned int    \\
    MPI_UNSIGNED_LONG   &   Unsigned long   \\
    MPI_FLOAT           &   Float           \\
    MPI_DOUBLE          &   Double          \\
    MPI_LONG_DOUBLE     &   Long Double     \\
    \multicolumn{2}{|c|}{...}               \\
    \hline
    \end{tabular}
    \end{center}
\end{table}
\end{frame}

\begin{frame}{Режимы операций передачи сообщений}
Режим MPI--коммуникаций определяет, при каких условиях операции передачи
завершаются.
\vfill
Режимы могут быть блокирующими (\abbr blocking) или неблокирующими (\abbr
non--blocking):
\begin{itemize}
    \item Blocking --- возврат из функции передачи сообщений только при
    завершении коммуникаций.
    \item Non--blocking --- немедленный возврат из функции, контроль завершения
    передачи лежит на программисте.
\end{itemize}
\end{frame}

\begin{frame}{Режимы передачи}
\begin{table}[htpb]
    \begin{center}
    \begin{tabularx}{\textwidth}{|l|X|}
    \hline
    Режим               &   Условие завершения                      \\
    \hline
    \hline
    Synchronous send    &   Завершение только после инициации приема \\
    \hline
    Standart send       &   Сообщение отослано (состояние приема неизвестно)    \\
    \hline
    Buffered send       &   Всегда завершается (кроме ошибочных передач), независимо от приема. \\
    \hline
    Ready send          &   Всегда завершается (кроме ошибочных передач), независимо от приема  \\
    \hline
    Receive             &   Завершается после приема                \\
    \hline
    \end{tabularx}
    \end{center}
\end{table}
\end{frame}

\begin{frame}{Функции передачи <<точка--точка>>}
\begin{table}[htpb]
    \begin{center}
    \begin{tabular}{|l|l|}
    \hline
    Режим               &   MPI функция \\
    \hline
    \hline
    Synchronous send    &   MPI_Ssend   \\
    \hline
    Standart send       &   MPI_Send    \\
    \hline
    Buffered send       &   MPI_Bsend   \\
    \hline
    Ready send          &   MPI_Rsend   \\
    \hline
    Receive             &   MPI_Recv    \\
    \hline
    \end{tabular}
    \end{center}
\end{table}
\end{frame}

\begin{frame}{Блокирующие и неблокирующие операции}
\texttt{MPI_\{S|B|R\}Send/Recv} --- блокирующиеся функции.
\begin{figure}
\centering
\begin{tikzpicture}
    \fill[red] (0,0) circle (0.15\textwidth);
    \fill[white] (-0.1\textwidth,-0.03\textwidth) rectangle (0.1\textwidth,0.03\textwidth);
\end{tikzpicture}
\end{figure}
\texttt{MPI_Isend/Irecv} --- не блокирующиеся варианты \texttt{MPI_Send/Recv}.
\end{frame}

\begin{frame}[fragile]{Send/Receive Overview}
\begin{lstlisting}
int MPI_Send(void *buf, int count, MPI_Datatype datatype, int dest, int tag, MPI_Comm comm);

int MPI_Recv(void *buf, int count, MPI_Datatype datatype, int source, int tag, MPI_Comm comm, MPI_Status status);
\end{lstlisting}
\end{frame}

\begin{frame}{MPI_Status \& tag}
\texttt{MPI_Status}:

\begin{itemize}
    \item Содержит:
    \begin{itemize}
        \item Source: \texttt{status.MPI_SOURCE},
        \item Tag:  \texttt{status.MPI_TAG},
        \item Count: \texttt{MPI_Get_count}.
    \end{itemize}
    \item \texttt{MPI_STATUS_IGNORE} --- специальное значение, которое уменьшает количество используемой памяти в том случае, если вы не собираетесь рассматривать это поле.
\end{itemize}
\vfill
tag:
\begin{itemize}
    \item Число для идентификации сообщения,
    \item \texttt{MPI_ANY_TAG}.
\end{itemize}
\end{frame}

\begin{frame}[fragile]{Полученное сообщение}
Может быть меньше, чем указано в MPI_Recv.
\vfill
\texttt{count} --- число реально полученных элементов.
\vfill
\begin{lstlisting}
int MPI_Get_count(MPI_Status *status, MPI_Datatype datatype, int *count);
\end{lstlisting}
\end{frame}

\section{Замер времени}

\begin{frame}[fragile]{Вычисление времени}
\begin{lstlisting}
double MPI_Wtime();
\end{lstlisting}

Возвращает время в секундах, прошедшее с  некоторого момента в прошлом (точки отсчета). Гарантируется, что эта точка отсчета не будет изменена в течение жизни процесса.

\begin{lstlisting}
double MPI_Wtick();
\end{lstlisting}

Возвращает разрешение таймера (минимальное значение кванта времени).

\begin{lstlisting}
int MPI_Barrier(MPI_Comm comm);
\end{lstlisting}

Блокирует работу вызвавшего ее процесса до тех пор, пока все другие процессы группы также не вызовут эту функцию.
\end{frame}

\begin{frame}[fragile]{Пример}
\begin{lstlisting}[basicstyle=\scriptsize]
MPI_Barrier(MPI_COMM_WORLD);
if (rank == 0)
    start_time = MPI_Wtime();

// Parallel code to measure here.

MPI_Barrier(MPI_COMM_WORLD);
if (rank == 0)
    printf("Execution time: %lf\n", MPI_Wtime() - start_time);
\end{lstlisting}
\end{frame}

\section{Компиляция и запуск}

\begin{frame}[fragile]{Компиляция и запуск}
Для компиляции используется \texttt{mpicc} вместо \texttt{gcc}.
\vfill
Запуск производится командой вида

\begin{lstlisting}
> mpirun –np <n_processes> <programm> <args>
\end{lstlisting}
\end{frame}

\section*{Конец}

\begin{frame}[allowframebreaks]{Рекомендуемая литература}
\begin{thebibliography}{99}
    \bibitem{} MPI Documents. \url{http://www.mpi-forum.org/docs/docs.html}
    \bibitem{} MPI standart. \url{http://www.mcs.anl.gov/research/projects/mpi}
    \bibitem{} \textit{Оленев~Н.Н.} Основы параллельного программирования в
    системе MPI. --- 2005 --- 81~с.
\end{thebibliography}
\end{frame}

\begin{frame}{Задания}
\begin{itemize}
    \item MPI ring,
    \item MPI ping-pong.
\end{itemize}
\end{frame}

\begin{frame}{На следующей лекции}
\begin{itemize}
    \item Общие понятия и определения,
    \item Состояние гонки,
    \item Синхронизация,
    \item Ошибки синхронизации.
\end{itemize}
\end{frame}

\begin{frame}

{\huge{Спасибо за внимание!}\par}

\vfill

\tiny{\textit{Замечание}: все торговые марки и логотипы, использованные в данном материале, являются собственностью их владельцев. Представленная здесь точка зрения отражает личное мнение автора, не выступающего от лица какой-либо организации.}

\end{frame}

\end{document}
